\documentclass[../../main.tex]{subfiles}
\begin{document}

\begin{center}
  {\fontsize{48pt}{58pt}\selectfont \textbf{1968}\normalsize\textbf{年度}}
\end{center}

\begin{table}[htb]
  \centering
  \caption{出題テーマと難易度}
  \begin{tabular}{|c|p{9cm}|c|}
    \hline
    \textbf{問題番号} & \textbf{テーマ} & \textbf{難易度} \\
    \hline
    第1問 & 整数問題（不等式） & \\
    \hline
    第2問 & 不等式（二次関数の判別式） & \\
    \hline
    第3問 & 複素数平面と領域図示 & \\
    \hline
    第4問 & 軌跡・領域の図示 & \\
    \hline
    第5問 & 極限（順列と対数） & \\
    \hline
    第6問 & 積分の最小値 & \\
    \hline
  \end{tabular}
\end{table}

\begin{table}[htb]
  \centering
  \caption{解答}
  \renewcommand{\arraystretch}{1.4}
  \begin{tabular}{|c|c|p{9cm}|}
    \hline
    \textbf{問題番号} & \textbf{小問} & \textbf{答え} \\
    \hline
    第1問 & - & $(a,b,c)=(4,3,2),(5,3,2),(6,3,2),(7,3,2)$ \\
    \hline
    第2問 & - & $-\dfrac{5}{3}\le\dfrac{a}{b}\le-\dfrac{3}{5}$ または $\dfrac{a}{b}=1$ \\
    \hline
    第3問 & - & 図示（$\cos\dfrac{\alpha-\beta}{2}\cos\dfrac{\alpha}{2}\cos\dfrac{\beta}{2}\le0$の領域） \\
    \hline
    第4問 & - & 図示（直線$l$の通らない領域，境界含まず） \\
    \hline
    第5問 & - & $\dfrac{4}{e}$ \\
    \hline
    第6問 & - & $\dfrac{\pi}{16}$（$f(x)=\sin^2x\cos^2x$） \\
    \hline
  \end{tabular}
\end{table}

\paragraph{試験形態}
6問．試験時間は要確認．

\end{document}
